\documentclass[12pt]{article}

% Layout.
\usepackage[top=1.2in, bottom=0.75in, left=1in, right=1in, headheight=1.0in, headsep=0pt]{geometry}

% Fonts.
\usepackage{mathptmx}
\usepackage[scaled=0.86]{helvet}
\renewcommand{\emph}[1]{\textsf{\textbf{#1}}}

% TiKZ.
\usepackage{tikz, pgfplots}
\usetikzlibrary{calc}
\pgfplotsset{my style/.append style={axis x line=middle, axis y line=middle, xlabel={$x$}, ylabel={$y$}}}
\pgfplotsset{compat=1.18}

% Misc packages.
\usepackage{amsmath,amssymb,latexsym}
\usepackage{graphicx}
\usepackage{array}
\usepackage{xcolor}
\usepackage{multicol}

% Commands to set various header/footer components.
\makeatletter
\def\doctitle#1{\gdef\@doctitle{#1}}
\doctitle{Use {\tt\textbackslash doctitle\{MY LABEL\}}.}
\def\docdate#1{\gdef\@docdate{#1}}
\docdate{Use {\tt\textbackslash docdate\{MY DATE\}}.}
\def\doccourse#1{\gdef\@doccourse{#1}}
\let\@doccourse\@empty
\def\docscoring#1{\gdef\@docscoring{#1}}
\let\@docscoring\@empty
\def\docversion#1{\gdef\@docversion{#1}}
\let\@docversion\@empty
\makeatother

% Headers and footers layout.
\makeatletter
\usepackage{fancyhdr}
\pagestyle{fancy}
\fancyhf{} % Clears all headers/footers.
\lhead{\emph{\@doctitle\hfill\@docdate} \medskip
\ifnum \value{page} > 1\relax\else\\
\emph{Name: \rule{3.5in}{1pt}\hfill \@docscoring}
\fi}

\rfoot{\emph{\@docversion}}
\lfoot{\emph{\@doccourse}}
\cfoot{\emph{\thepage}}
\renewcommand{\headrulewidth}{0pt}%
\makeatother

% Paragraph spacing
\parindent 0pt
\parskip 6pt plus 1pt

% A problem is a section-like command. Use \problem{5} for a problem worth 5 points.
\newcounter{probcount}
\newcounter{subprobcount}
\setcounter{probcount}{0}

\newcommand{\problem}[1]{%
\par
\addvspace{4pt}%
\setcounter{subprobcount}{0}%
\stepcounter{probcount}%
\makebox[0pt][r]{\emph{\arabic{probcount}.}\hskip1ex}\emph{[#1 points]}\hskip1ex}

\newcommand{\thesubproblem}{\emph{\alph{subprobcount}.}}

% like \problem but with name
\newcommand{\nameprob}[2]{%
\par
\addvspace{4pt}%
\setcounter{subprobcount}{0}%
\stepcounter{probcount}%
\makebox[0pt][r]{\emph{#1.}\hskip1ex}\emph{[#2 points]}\hskip1ex}

% Subproblems are an enumerate-like environment with a consistent
% numbering scheme.  Use \begin{subproblems}\item...\item...\end{subproblems}
\newenvironment{subproblems}{%
\begin{enumerate}%
\setcounter{enumi}{\value{subprobcount}}%
\renewcommand{\theenumi}{\emph{\alph{enumi}}}}%
{\setcounter{subprobcount}{\value{enumi}}\end{enumerate}}

% Blanks for answers in normal and math mode.
\newcommand{\blank}[1]{\rule{#1}{0.75pt}}
\newcommand{\mblank}[1]{\underline{\hspace{#1}}}
\def\emptybox(#1,#2){\framebox{\parbox[c][#2]{#1}{\rule{0pt}{0pt}}}}

% Misc.
\renewcommand{\d}{\displaystyle}
\newcommand{\ds}{\displaystyle}


\doctitle{Math 252: Quiz 7 }
\docdate{23 October 2025}
\doccourse{}
\docversion{v1}
\docscoring{\fbox{{\LARGE \strut}\blank{0.8in} / 25}}

\begin{document}
30 minutes maximum.  No aids (book, calculator, etc.) are permitted.  Show all work and use proper notation for full credit.  Answers should be in reasonably-simplified form.  25 points possible.

\begin{enumerate}
\item (5 points) Consider the \textbf{telescoping} series $\ds \sum_{n=1}^{\infty} \left( \frac{1}{n}- \frac{1}{n+2}\right).$
	\begin{enumerate}
%	\item Find $S_3,$ the third partial sum of the series.  (You do not need to simplify your numerical answer.)
%	\vfill
	\item Find a simplified formula for $S_k$, the $k$th partial sum of the series.
	\vfill
	\item Use your answer from part (b) to determine the sum of the series.\\
	\vfill
	\end{enumerate}
\item (8 points) For each \textbf{geometric} series below, determine whether the series converges and \textbf{justify} your conclusion. \textbf{If} the series converges, determine the sum.
	\begin{enumerate}
	\item $\ds \sum_{n=1}^{\infty} \frac{(-4)^{n-1}}{\pi 5^{n-1}}$
	\vfill
	\item $\ds \sum_{n=1}^{\infty} \frac{2^{2n}}{3^n}$
	\vfill
	\end{enumerate}
\newpage
\item (8 points) For each  series below, determine whether the series converges and \textbf{justify} your conclusion. Your justification must include the test or series you are using to draw your conclusion.
	\begin{enumerate}
	\item $\ds \sum_{n=1}^{\infty} \frac{n}{10n+17}$
	\vfill
	\item $\ds \sum_{n=1}^{\infty} \frac{1}{n^{\sqrt{5}}}$
	\vfill
	\end{enumerate}
\item (4 points) Apply the Integral Test to determine if the series $\ds \sum_{n=2}^{\infty} \frac{1}{n \sqrt{\ln (n)}}$ converges or diverges. (You do not have to show that the Integral Test applies; it does.)
\vfill
\vfill
\end{enumerate}
\end{document}

